\documentclass[11pt,letterpaper]{article}
\usepackage[margin=1in,headheight=15pt]{geometry}
\usepackage[T1]{fontenc}
\usepackage{amsmath,amsthm,mathtools}
\usepackage{newtxtext,newtxmath}
\usepackage{microtype}
\usepackage{booktabs,array,longtable}
\usepackage[dvipsnames]{xcolor}
\usepackage{enumitem}
\usepackage[most]{tcolorbox}
\usepackage{fancyhdr}
\usepackage{listings}
\usepackage{hyperref}
\usepackage{bookmark}
\definecolor{Ink}{HTML}{17324D}
\definecolor{Accent}{HTML}{176B78}
\definecolor{Pale}{HTML}{F1F6F8}
\hypersetup{colorlinks=true,linkcolor=Ink,citecolor=Accent,urlcolor=Accent,
 pdftitle={Finite Hankel Certificates for Three Somos-8 Conjectures},
 pdfauthor={Research report prepared with ChatGPT},
 pdfsubject={Proofs of the Somos-8 assertions in Barry's Conjectures 11--13}}
\pagestyle{fancy}
\fancyhf{}
\fancyhead[L]{\small\scshape Finite Hankel certificates}
\fancyhead[R]{\small Somos--8 conjectures}
\fancyfoot[C]{\thepage}
\renewcommand{\headrulewidth}{0.4pt}
\setlength{\parskip}{4pt}
\setlength{\parindent}{1.2em}
\allowdisplaybreaks[2]
\numberwithin{equation}{section}
\newtheorem{theorem}{Theorem}[section]
\newtheorem{lemma}[theorem]{Lemma}
\newtheorem{proposition}[theorem]{Proposition}
\newtheorem{corollary}[theorem]{Corollary}
\theoremstyle{definition}
\newtheorem{definition}[theorem]{Definition}
\newtheorem{remark}[theorem]{Remark}
\newcommand{\Z}{\mathbb Z}
\newcommand{\Q}{\mathbb Q}
\newcommand{\C}{\mathbb C}
\newcommand{\coef}[2]{[#1^{#2}]}
\newcommand{\calH}{\mathcal H}
\newcommand{\calV}{\mathcal V}
\newcommand{\Res}{\mathcal R}
\newcommand{\Bcal}{\mathcal B}
\newtcolorbox{resultbox}[1][]{enhanced,breakable,colback=Pale,colframe=Accent,
 boxrule=0.65pt,arc=1.5mm,left=3mm,right=3mm,top=2mm,bottom=2mm,#1}
\lstset{basicstyle=\ttfamily\small,breaklines=true,columns=fullflexible,
 keepspaces=true,showstringspaces=false,frame=single,rulecolor=\color{Accent},
 backgroundcolor=\color{Pale}}

\begin{document}
\begin{titlepage}
\color{Ink}
\noindent\textsc{Integer sequences \enspace / \enspace generating functions \enspace / \enspace exact proof}
\vspace{1.0in}
\begin{flushleft}
{\Huge\bfseries Finite Hankel Certificates\par}
\vspace{0.12in}
{\Huge\bfseries for Three Somos--8\par}
\vspace{0.12in}
{\Huge\bfseries Conjectures\par}
\vspace{0.35in}
{\Large The Somos--8 assertions in Barry's Conjectures 11--13,\par}
\vspace{0.07in}
{\Large singular specializations, and coefficient asymptotics\par}
\end{flushleft}
\vspace{0.45in}
\noindent\textbf{Prepared for Vladimir Reshetnikov}\par
\noindent 19 September 2026
\vspace{0.35in}
\begin{resultbox}[title={\bfseries Result and proof status}]
The three parameterized bilinear Hankel identities are established by combining
Hone's published genus-two theorem with a specialization argument and twelve
exact polynomial identities. The local certificates are reproduced by a
standard-library Python verifier and independently checked with SymPy.

The proof does not assume the conjectured recurrences to generate its data.
It remains valid in denominator-cleared form at singular parameter values.
For Conjecture 13, the defining continued fraction is used; a sign error in the
source's separate Catalan display is corrected explicitly.

No earlier resolution of these three Somos--8 assertions was located in the
focused literature search. This is not an exhaustive certification of novelty.
\end{resultbox}
\vfill
\noindent\small Research article and reproducibility package. The genus-two existence theorem
is an external mathematical input, not a theorem reproved in this report.
\end{titlepage}

\begin{abstract}
Barry proposed three one-parameter families of algebraic generating functions
whose Hankel transforms were conjectured to satisfy explicit Somos--8
recurrences. We prove the three bilinear assertions, namely the Somos--8
parts of Conjectures 11, 12, and 13 in arXiv:2211.12637. The argument makes a
finite computation logically sufficient: Hone's genus-two Hankel theorem
implies that an infinite matrix of quadratic products has rank at most four;
a binomial translation puts each proposed family into the corresponding
universal moment class; and four symbolic rows determine the recurrence.
A polynomial-specialization argument is essential because the proposed
families lie on a degenerate parameter locus of the usual continued-fraction
coordinates. We give denominator-free formulas, exact nonzero rank witnesses,
and executable integer-polynomial certificates. For the second family we
also derive two Somos--6 specializations, a positive finite coefficient
formula with a tree interpretation, and an all-orders dominant-singularity
expansion of its generating-function coefficients. These asymptotics concern
the moments, not their Hankel determinants.
\end{abstract}

\setcounter{tocdepth}{1}
\tableofcontents
\newpage

\section{The selected problem and its provenance}
\label{sec:provenance}

A useful research target should have an unambiguous published statement, enough
structure to support an all-index argument, and a feasible independent
verification path. The Somos--8 conjectures in Barry's 2022 preprint
\cite{Barry} meet these requirements. They concern an algebraic power series
\(g(x)=\sum_{n\geq0}a_nx^n\) and its Hankel transform
\begin{equation}
 h_n=\det(a_{i+j})_{0\leq i,j\leq n},\qquad n\geq0.
 \label{eq:hankel}
\end{equation}
The striking assertion is not integrality of \(h_n\), which follows directly
when the \(a_n\) are integers. It is that the determinants obey a fixed
short bilinear recurrence with explicitly conjectured rational coefficients.

Our main target is Conjecture 12. Its generating function is the unique formal
power-series solution with constant term one of
\begin{equation}
 g(x)=\frac{1}{1-x-\dfrac{x^2}{1-rx}-x^3g(x)}.
 \label{eq:c12cf}
\end{equation}
The proposed Hankel recurrence has order eight. The proof developed for this
family also proves the adjacent Conjectures 11 and 13 with no new structural
input, so all three are treated together.

There are important distinctions in the literature. Wang and Zhang
\cite{WangZhang} prove the four Somos--4 conjectures in Barry's paper, not the
three Somos--8 assertions studied here. Hone's earlier work \cite{Hone}
provides the general genus-two theorem that makes our proof possible, but its
statement does not list Barry's later one-parameter coefficient formulas.
The July 2026 paper of Liu, Wang, and Zhang \cite{LiuWangZhang}, despite its
similar title, concerns four conjectures about a \emph{central transform}.
The September 2026 note of Scheuerle \cite{Scheuerle} studies inverse
realization of prescribed Somos--4 Hankel data. These are different problems.

The search leading to this report checked the original preprint, the cited
Somos--4 resolution, Hone's theorem and its moment parameterization, and
recent papers surfaced by searches for Barry's Hankel conjectures. No prior
proof of the three precise Somos--8 formulas was located. This supports
selecting them as conjectural targets, but it does not prove that no such
proof exists elsewhere. The mathematical claims below are proofs of the
stated identities; historical priority is a separate question.

\subsection{What is established, and what is imported}

The only substantial structural input is Hone's genus-two theorem
\cite[Theorem 5.4]{Hone}. Its use is explicit. We prove the translation,
the polynomial-specialization step, the finite-certification lemma, the
parameter substitutions, the exceptional-parameter consequences, and the
coefficient formulas and asymptotics. The finite identities needed to
identify the recurrence are evaluated exactly in \(\Z[r]\), using two
independent determinant implementations.

This is therefore not merely a numerical confirmation, and it is not a
self-contained rederivation of genus-two integrability. It is a proof by
reduction to a published theorem, with all family-specific algebra made
reproducible. We do not claim a proof of the other conjectures in Barry's
paper, nor a classification of all algebraic generating functions with
Somos Hankel transforms.

\section{Precise statements of the three identities}
\label{sec:main}

\subsection{Uniform notation}

Let \(r\) be an indeterminate and set \(B(x)=1-rx\). Define
\begin{align}
 A_{11}(x)&=1-(r+1)x-x^2+rx^3,\label{eq:A11}\\
 A_{12}(x)&=1-(r+1)x+(r-1)x^2,\label{eq:A12}\\
 A_{13}(x)&=1-(r+1)x+(r-2)x^2+rx^3.\label{eq:A13}
\end{align}
For \(k\in\{11,12,13\}\), define \(g^{(k)}(x)\in\Z[r][[x]]\) by
\begin{equation}
 A_k(x)g^{(k)}(x)=B(x)+x^3B(x)g^{(k)}(x)^2,
 \qquad g^{(k)}(0)=1.
 \label{eq:uniform}
\end{equation}
Write its coefficients as \(a_n^{(k)}\), and use \eqref{eq:hankel} for
\(h_n^{(k)}\). We also define the empty determinant
\begin{equation}
 h_{-1}^{(k)}=1.
 \label{eq:empty}
\end{equation}
The extra index allows the bilinear identity to start at \(n=7\), one step
before the customary recurrence written solely with nonnegative indices.

Dividing \eqref{eq:uniform} by \(B(x)g^{(k)}(x)\) yields respectively
\begin{align}
 g^{(11)}(x)&=\frac1{1-\dfrac{x}{1-rx}-x^2-x^3g^{(11)}(x)},\label{eq:CF11}\\
 g^{(12)}(x)&=\frac1{1-x-\dfrac{x^2}{1-rx}-x^3g^{(12)}(x)},\label{eq:CF12}\\
 g^{(13)}(x)&=\frac1{1-\dfrac{x(1-(r-1)x)}{1-rx}-x^2-x^3g^{(13)}(x)}.
 \label{eq:CF13}
\end{align}
These are the defining continued fractions in Barry's Conjectures 11--13.

\begin{remark}[A source correction, not a counterexample]
In the separate Catalan generating-function display for Conjecture 13,
Barry prints \(1+(r+1)x+(r-2)x^2+rx^3\). The defining continued fraction
\eqref{eq:CF13} instead gives \eqref{eq:A13}, with a minus sign before
\((r+1)x\). This follows by multiplying its denominator by \(1-rx\).
We use the continued-fraction definition throughout. The claimed Somos
parameters require no change. Our result establishes that substantive
recurrence claim, not the inconsistent pair of generating-function displays
read literally.
\end{remark}

\subsection{Denominator-free coefficient tables}

The following polynomials specify the answer. The subscript on \(P_{k,j}\)
is the family number followed by the position of the quadratic product.
Using a common denominator avoids spurious exclusions at its roots.

For the family \(k=12\), the simplest of the three, let
\begin{align}
 D_{12}&=r^2-4r+3=(r-1)(r-3),\label{eq:D12}\\
 P_{12,1}&=r^4-11r^3+26r^2-16r-5,\label{eq:P121}\\
 P_{12,2}&=2r^5-19r^4+40r^3-13r^2-5r,\label{eq:P122}\\
 P_{12,3}&=(r-1)(3r^6-12r^5+15r^4+25r^3-62r^2-36r-5),\label{eq:P123}\\
 P_{12,4}&=r^9-8r^8+26r^7-43r^6+40r^5-17r^4\notag\\
 &\hspace{1.8em}{}-23r^3+27r^2+19r+3.\label{eq:P124}
\end{align}
For \(k=11\), put
\begin{align}
 D_{11}&=r^4-2r^3+8r^2+2r-9
       =(r^2-1)(r^2-2r+9),\label{eq:D11}\\
 P_{11,1}&=r^8-8r^7+21r^6-40r^5+35r^4-24r^3+71r^2+8r,\\
 P_{11,2}&=8(r^9-6r^8+17r^7-30r^6+15r^5-14r^4-r^3-14r^2),\\
 P_{11,3}&=8(r-1)\bigl(r^{10}-2r^8+29r^7-32r^6+39r^5\notag\\
 &\hspace{6.7em}{}+18r^4+11r^3-r^2+r\bigr),\\
 P_{11,4}&=2r^{13}-13r^{12}+48r^{11}-85r^{10}+83r^9-11r^8
             +124r^7\notag\\
 &\hspace{1.8em}{}-454r^6+364r^5-263r^4-84r^3-189r^2-25r-9.
 \label{eq:P11last}
\end{align}
Finally, for \(k=13\), put
\begin{align}
 D_{13}&=2(r^3-3r^2-5r+7)=2(r-1)(r^2-2r-7),\label{eq:D13}\\
 P_{13,1}&=-(r^7-8r^6+25r^5-20r^4-37r^3+75r+28),\\
 P_{13,2}&=(r+1)\bigl(r^8-11r^7+47r^6-83r^5+17r^4\notag\\
 &\hspace{6.2em}{}+71r^3+45r^2-169r+210\bigr),\\
 P_{13,3}&=(r^2-1)\bigl(3r^8-29r^7+115r^6-225r^5+181r^4\notag\\
 &\hspace{6.7em}{}+105r^3-255r^2-235r+84\bigr),\\
 P_{13,4}&=-\bigl(r^{10}-17r^9+96r^8-212r^7+54r^6+594r^5\notag\\
 &\hspace{3.5em}{}-796r^4-36r^3+721r^2-329r-588\bigr).
 \label{eq:P13last}
\end{align}

\begin{theorem}[Three Somos--8 identities]
\label{thm:main}
For every \(k\in\{11,12,13\}\) and every integer \(n\geq7\), the identity
\begin{equation}
 \boxed{\quad
 D_k(r)h_n^{(k)}h_{n-8}^{(k)}
 =\sum_{j=1}^4P_{k,j}(r)h_{n-j}^{(k)}h_{n-8+j}^{(k)}
 \quad}
 \label{eq:main}
\end{equation}
holds in \(\Z[r]\). In particular, for a complex parameter \(r\) with
\(D_k(r)\ne0\), these Hankel determinants satisfy the Somos--8 bilinear
relation with coefficients
\begin{equation}
 (\alpha,\beta,\gamma,\delta)
 =D_k(r)^{-1}(P_{k,1},P_{k,2},P_{k,3},P_{k,4}).
 \label{eq:normalized}
\end{equation}
These are precisely the Somos--8 coefficient formulas in Barry's
Conjectures 11, 12, and 13, interpreted using their defining continued
fractions. At roots of \(D_k\), the polynomial relation \eqref{eq:main}
still holds.
\end{theorem}

The proof is given in Sections~\ref{sec:rank}--\ref{sec:certificate}.
For \(k=11\), the factor \(r-1\) in \(P_{11,3}\) explains why the source
uses a smaller denominator for \(\gamma\); for \(k=12\), the analogous
factor accounts for its denominator \(r-3\). Thus no coefficients have
been fitted approximately or altered to obtain the theorem.

\begin{corollary}[Integrality and arbitrary specializations]
\label{cor:integral}
For every integer \(r\), all moments and Hankel determinants are integers.
More generally, \eqref{eq:main} specializes to every commutative ring with
identity by sending \(r\) to any element of that ring.
\end{corollary}
\begin{proof}
The coefficient recursion from \eqref{eq:uniform} uses only addition and
multiplication in \(\Z[r]\); see \eqref{eq:momentrec} below. Determinants
therefore also belong to \(\Z[r]\). Apply a ring homomorphism to the
polynomial identity in Theorem~\ref{thm:main}.
\end{proof}
This last statement does not apply a genus-two analytic theorem directly in
positive characteristic. It specializes an integral polynomial identity
already proved in characteristic zero.

\section{Formal power series and exact moment generation}
\label{sec:moments}

Write \(A_k(x)=1+A_1x+A_2x^2+A_3x^3\), with \(A_3=0\) when \(k=12\).
In all three cases \(A_1=-(r+1)\). Set \(a_n=0\) for negative \(n\), and
\begin{equation}
 c_m=\begin{cases}\displaystyle\sum_{i=0}^m a_i a_{m-i},&m\geq0,\\
                  0,&m<0.
      \end{cases}
\end{equation}
Equating coefficients in \eqref{eq:uniform} gives \(a_0=1\) and
\begin{equation}
 a_n=-A_1a_{n-1}-A_2a_{n-2}-A_3a_{n-3}
      +c_{n-3}-rc_{n-4}-r\,\mathbf1_{n=1},\qquad n\geq1.
 \label{eq:momentrec}
\end{equation}
Every term on the right is already known at stage \(n\). This proves both
existence and uniqueness of the formal series, as well as the asserted
integral-polynomial coefficient property. The first coefficients are
\begin{equation}
 a_0=1,\qquad a_1=1,\qquad
 a_2=\begin{cases}r+2,&k=11,\\2,&k=12,\\3,&k=13.\end{cases}
 \label{eq:firstmoments}
\end{equation}

If \(C(z)=\sum_{m\geq0}C_mz^m\) is the Catalan series, characterized by
\(C(z)=1+zC(z)^2\), then
\begin{equation}
 g^{(k)}(x)=\frac{B(x)}{A_k(x)}
 C\!\left(\frac{x^3B(x)^2}{A_k(x)^2}\right).
 \label{eq:catalan}
\end{equation}
This follows by substitution, or from the branch
\begin{equation}
 g^{(k)}(x)=\frac{A_k(x)-\sqrt{A_k(x)^2-4x^3B(x)^2}}{2x^3B(x)},
 \label{eq:root}
\end{equation}
where the square root has constant term one. The apparent singularity at
\(x=0\) in \eqref{eq:root} is removable as a formal series.

It is essential that our computation proceeds from \eqref{eq:momentrec}
and then evaluates determinants. Generating \(h_n\) from the proposed
Somos recurrence and comparing those values with the same recurrence
would be circular. No such step is used in the verifier.

For orientation, the first few Hankel polynomials for \(k=12\) are
\begin{align}
 h_{-1}=h_0=h_1&=1,\qquad h_2=-2r,\notag\\
 h_3&=-r^4+2r^3-r^2-4r-1,\notag\\
 h_4&=-r^6+4r^5-5r^4+r^3-6r^2-5r-1,\notag\\
 h_5&=r^7-5r^6+r^5+12r^4+12r^3-7r^2-6r-1.
 \label{eq:initialpolys}
\end{align}
The relatively small degrees of these polynomials make exact family-level
certification much less expensive than a general seven-parameter calculation.

\section{The structural input: a rank-four theorem}
\label{sec:rank}

\subsection{A universal genus-two moment class}

Introduce seven algebraically independent parameters
\(u,v,b_2,b_3,c_2,c_3,c_4\), and moments
\(s_0=1,s_1=u,s_2=v\). For \(m\geq3\), define
\begin{align}
 s_m={}&b_2s_{m-2}+b_3s_{m-3}
       +c_2\sum_{i=0}^{m-2}s_is_{m-2-i}\notag\\
     &+c_3\sum_{i=0}^{m-3}s_is_{m-3-i}
       +c_4\sum_{i=0}^{m-4}s_is_{m-4-i},
 \label{eq:universal}
\end{align}
with an empty sum equal to zero. Let
\begin{equation}
 H_N=\det(s_{i+j})_{0\leq i,j<N},\qquad H_0=1.
 \label{eq:HN}
\end{equation}
Notice the size convention: \(H_N\) is an \(N\)-by-\(N\) determinant,
whereas \(h_n\) is an \((n+1)\)-by-\((n+1)\) determinant.

\begin{theorem}[Hone's genus-two theorem, the form used here]
\label{thm:hone}
For the generic moment sequence \eqref{eq:universal}, the determinants
\eqref{eq:HN} satisfy a nontrivial relation
\begin{equation}
 q_0H_{N+8}H_N+q_1H_{N+7}H_{N+1}
 +q_2H_{N+6}H_{N+2}+q_3H_{N+5}H_{N+3}
 +q_4H_{N+4}^2=0,
 \label{eq:Hone}
\end{equation}
where the coefficients are independent of \(N\).
\end{theorem}

This is the input from \cite[Theorem 5.4, equation (4.22), and the
seven-parameter formulation in its proof]{Hone}. We require existence and
index-independence of the relation, not the large general formulas for the
\(q_j\). In the generic setting, the relation holds for \(N\geq0\).

For clarity, the link from \eqref{eq:universal} to Hone's curve coordinates
can be exhibited directly. Put
\begin{equation}
 G(X)=\sum_{m\geq0}s_mX^{-m-1},
\end{equation}
and define
\begin{align}
 P(X)&=X^3-b_2X-b_3,\notag\\
 S(X)&=c_2X^2+c_3X+c_4,\notag\\
 T(X)&=X^2+uX+v-b_2-c_2.
 \label{eq:univPST}
\end{align}
The coefficient recursion is equivalent to
\begin{equation}
 P(X)G(X)-S(X)G(X)^2=T(X).
 \label{eq:PGST}
\end{equation}
Indeed, its polynomial part specifies \(s_0,s_1,s_2\), while each negative
power of \(X\) gives \eqref{eq:universal}. Setting
\begin{equation}
 Y=P-2SG
\end{equation}
yields
\begin{equation}
 Y^2=P^2-4ST.
 \label{eq:sextic}
\end{equation}
Generically this is a nonsingular sextic model of a genus-two curve. Moreover,
\(Q_0=-2S\) and \(Q_1=2T\) put \eqref{eq:PGST} into Hone's equation
\(PG+Q_0G^2/2=Q_1/2\). The polynomial part of the square root at infinity is
\begin{equation}
 X^3-(b_2+2c_2)X-(b_3+2c_2u+2c_3).
\end{equation}
Thus the seven parameters in \eqref{eq:universal} are the generic moment
parameters of the cited theorem, not an additional unproved generalization
of it.

\subsection{Why specialization needs an argument}

All moments in \eqref{eq:universal}, and hence all \(H_N\), are polynomials
in the seven parameters. Consider the infinite matrix whose row indexed by
\(N\geq0\) is
\begin{equation}
 (H_{N+8}H_N,\ H_{N+7}H_{N+1},\ H_{N+6}H_{N+2},\
  H_{N+5}H_{N+3},\ H_{N+4}^2).
 \label{eq:universalmatrix}
\end{equation}
Theorem~\ref{thm:hone} implies that its five columns are linearly dependent
over the generic parameter field.

\begin{lemma}[Polynomial rank specialization]
\label{lem:specialization}
Every five-by-five minor of \eqref{eq:universalmatrix}, using any five row
indices, is the zero polynomial. Consequently, after any specialization of
the seven parameters into a characteristic-zero field, the specialized
matrix has rank at most four.
\end{lemma}
\begin{proof}
Fix any five row indices. Their minor is a polynomial in the seven
parameters. It is zero in the generic field by Theorem~\ref{thm:hone}, so
it is the zero polynomial. Equivalently, it vanishes on the generic set
where the theorem applies and hence identically. Substitution in a zero
polynomial preserves zero. Since the chosen rows were arbitrary, all
five-by-five minors vanish after specialization.
\end{proof}

The lemma concerns \emph{arbitrary} row selections, not merely five
consecutive rows. The existence of a single index-independent relation in
Hone's theorem is what supplies this stronger conclusion. Also, the rank
cannot rise when specializing polynomial data, although it can drop.

This step matters here because the desired families will have \(c_2=0\).
In the direct continued-fraction chart, \(Q_0\) then loses degree and some
coordinate formulas divide by \(c_2\). Substituting zero into those rational
formulas would be invalid. Specializing the polynomial minors instead is
valid and is sufficient for the proof.

\section{Binomial translation and the three parameter substitutions}
\label{sec:translation}

\subsection{Hankel invariance under translating moments}

\begin{lemma}[Binomial translation]
\label{lem:binomial}
For any sequence \((a_n)\) and scalar \(c\), define
\begin{equation}
 s_n=\sum_{j=0}^n\binom nj(-c)^{n-j}a_j.
 \label{eq:binomial}
\end{equation}
Then for every \(N\geq0\),
\begin{equation}
 \det(s_{i+j})_{0\leq i,j<N}
 =\det(a_{i+j})_{0\leq i,j<N}.
 \label{eq:invariance}
\end{equation}
In Laurent-series notation, if \(F(X)=\sum_{n\geq0}a_nX^{-n-1}\), then
\begin{equation}
 F(X+c)=\sum_{n\geq0}s_nX^{-n-1}.
 \label{eq:Laurenttranslate}
\end{equation}
\end{lemma}
\begin{proof}
Let \(\mathscr L\) be the linear functional on polynomials with
\(\mathscr L(X^n)=a_n\). Then \(s_n=\mathscr L((X-c)^n)\). The change of
basis from \(1,X,\ldots,X^{N-1}\) to
\(1,X-c,\ldots,(X-c)^{N-1}\) is lower triangular with diagonal entries one.
If its matrix is \(U\), the new moment matrix is
\(U(a_{i+j})U^{\mathsf T}\), so its determinant is unchanged.
For the Laurent identity, expand each \((X+c)^{-j-1}\) using the binomial
series and collect the coefficient of \(X^{-n-1}\); it is exactly
\eqref{eq:binomial}.
\end{proof}

\subsection{Removing the quadratic term in the cubic}

For a fixed family, suppress its superscript. Put
\begin{equation}
 F(X)=X^{-1}g(X^{-1}),\qquad
 \mathcal A(X)=X^3A(X^{-1})=X^3-(r+1)X^2+A_2X+A_3.
\end{equation}
Substituting into \eqref{eq:uniform} and multiplying by \(X^2\) gives
\begin{equation}
 \mathcal A(X)F(X)-(X-r)F(X)^2=X(X-r).
 \label{eq:Laurentquadratic}
\end{equation}
The powers of \(X\) in this step are important: \(g(X^{-1})=XF(X)\),
not \(F(X)\).

Set
\begin{equation}
 c=\frac{r+1}{3},\qquad
 p=A_2-3c^2,\qquad q=A_3+A_2c-2c^3,
 \label{eq:cpq}
\end{equation}
and let \(G(X)=F(X+c)\). Since
\(\mathcal A(X+c)=X^3+pX+q\), equation \eqref{eq:Laurentquadratic} becomes
\begin{equation}
 (X^3+pX+q)G-(X+c-r)G^2=(X+c)(X+c-r).
 \label{eq:translatedquadratic}
\end{equation}
Its coefficients at negative powers of \(X\) give, for \(m\geq3\),
\begin{equation}
 s_m=-ps_{m-2}-qs_{m-3}
       +\sum_{i=0}^{m-3}s_is_{m-3-i}
       +(c-r)\sum_{i=0}^{m-4}s_is_{m-4-i}.
 \label{eq:specialmoment}
\end{equation}
The initial data are
\begin{equation}
 s_0=1,\qquad s_1=1-c,\qquad s_2=a_2-2c+c^2.
 \label{eq:specialinitial}
\end{equation}
Thus the required substitution into \eqref{eq:universal} is
\begin{equation}
 \boxed{\quad
 (u,v,b_2,b_3,c_2,c_3,c_4)
 =(1-c,\ a_2-2c+c^2,\ -p,\ -q,\ 0,\ 1,\ c-r).
 \quad}
 \label{eq:sevenparameters}
\end{equation}
The only family-dependent input is
\begin{center}
\begin{tabular}{c r r r}
\toprule
Family & \(A_2\) & \(A_3\) & \(a_2\)\\
\midrule
11 & \(-1\) & \(r\) & \(r+2\)\\
12 & \(r-1\) & \(0\) & \(2\)\\
13 & \(r-2\) & \(r\) & \(3\)\\
\bottomrule
\end{tabular}
\end{center}

By Lemma~\ref{lem:binomial}, the translated sequence has the same Hankel
determinants. Lemma~\ref{lem:specialization}, applied over \(\Q(r)\),
therefore proves that the infinite matrix with rows
\begin{equation}
 \bigl(h_nh_{n-8},\ h_{n-1}h_{n-7},\ h_{n-2}h_{n-6},\
 h_{n-3}h_{n-5},\ h_{n-4}^2\bigr),\qquad n\geq7,
 \label{eq:familymatrix}
\end{equation}
has rank at most four for each of the three families. We have established
an all-index structural bound before checking any of the conjectured
coefficient formulas.

\section{Why four exact rows complete the proof}
\label{sec:finite}

\begin{lemma}[Finite certification of a fixed recurrence]
\label{lem:finite}
Let \(K\) be a field and \(h_{-1},h_0,\ldots\) a sequence in \(K\).
For \(n\geq7\), set
\begin{equation}
 y_n=h_nh_{n-8},\qquad
 v_n=(h_{n-1}h_{n-7},\ h_{n-2}h_{n-6},\ h_{n-3}h_{n-5},\ h_{n-4}^2).
\end{equation}
Assume that the matrix with rows \((y_n,v_n)\) has rank at most four.
Choose four row indices \(n_1,\ldots,n_4\), and suppose that the matrix
\(M=(v_{n_i})_{i=1}^4\) is invertible. If \(D,P_1,\ldots,P_4\in K\)
satisfy
\begin{equation}
 Dy_{n_i}=v_{n_i}\begin{pmatrix}P_1&P_2&P_3&P_4\end{pmatrix}^{\mathsf T}
 \qquad(1\leq i\leq4),
 \label{eq:fourchecks}
\end{equation}
then the same identity holds for every \(n\geq7\).
\end{lemma}
\begin{proof}
Take the five-by-five matrix with columns \((v,y)\), using the four chosen
rows and an arbitrary fifth row \(n\). Its determinant is zero by the rank
hypothesis. Replace its last column by \(D\) times that column minus
\(\sum_{j=1}^4P_j\) times the corresponding first four columns.
The new determinant is still zero. By \eqref{eq:fourchecks}, its first
four entries in the last column are zero. Expanding on that column yields
\begin{equation}
 \det(M)\left(Dy_n-v_nP\right)=0.
\end{equation}
Invertibility of \(M\) gives the conclusion. A repeated fifth row causes
no difficulty; the desired identity for the original four rows was already
assumed.
\end{proof}

The rank hypothesis is indispensable. For an arbitrary sequence, four or
four million initial recurrence checks do not prove the next check. Here
the rank theorem bounds the entire infinite family of possible quadratic
rows in advance. The finite computation only identifies the one relation
already forced to exist.

\section{Exact certificates and completion of the main proof}
\label{sec:certificate}

We use the four rows \(n=7,8,9,10\). They require only
\(h_{-1},h_0,\ldots,h_{10}\), and thus only
\(a_0,a_1,\ldots,a_{20}\). For each family, \eqref{eq:momentrec}
computes these as integer polynomials, and direct determinants give the
needed Hankel polynomials.

Define the residual polynomials
\begin{equation}
 E_{k,n}(r)=D_kh_nh_{n-8}
 -\sum_{j=1}^4P_{k,j}h_{n-j}h_{n-8+j}.
 \label{eq:residual}
\end{equation}
The exact computations yield
\begin{center}
\begin{tabular}{c c c c c}
\toprule
Family & \(E_{k,7}\) & \(E_{k,8}\) & \(E_{k,9}\) & \(E_{k,10}\)\\
\midrule
11 & \(0\) & \(0\) & \(0\) & \(0\)\\
12 & \(0\) & \(0\) & \(0\) & \(0\)\\
13 & \(0\) & \(0\) & \(0\) & \(0\)\\
\bottomrule
\end{tabular}
\end{center}
Here each zero means the \emph{zero polynomial in \(\Z[r]\)}, not a
floating-point value and not an interpolation guess.

The polynomial degrees involved are modest:
\begin{center}
\small
\begin{tabular}{c rrrrrrrrrrr}
\toprule
 & \multicolumn{11}{c}{Degree of \(h_n(r)\)}\\
Family / \(n\) &0&1&2&3&4&5&6&7&8&9&10\\
\midrule
11&0&1&3&5&7&11&15&19&24&30&36\\
12&0&0&1&4&6&7&10&15&19&22&27\\
13&0&0&2&3&6&8&11&15&18&23&27\\
\bottomrule
\end{tabular}
\end{center}
The actual coefficient arrays, including every polynomial needed to rerun
the verification, are recorded in \texttt{data/certificates.json}.

To prove that the four product rows are independent over \(\Q(r)\), one
nonzero specialization suffices. At \(r=2\), the respective four-by-four
determinants are
\begin{align}
 \det M_{11}(2)&=844298284526087295172008083616,\notag\\
 \det M_{12}(2)&=-1967013462528657296,\notag\\
 \det M_{13}(2)&=-83222140602851259464.
 \label{eq:witnesses}
\end{align}
In particular none of the three determinant polynomials is identically zero.
For the central family, the entire witness is small enough to display:
\begin{equation}
 M_{12}(2)=
 \begin{pmatrix}
 2512&87&172&169\\
 22041&-10048&-1131&1849\\
 -719780&-286533&-108016&7569\\
 89140155&-7737635&1917567&6310144
 \end{pmatrix}.
 \label{eq:M12}
\end{equation}
The other two matrices are included in Appendix~\ref{app:witness} and the
JSON certificate.

\begin{proof}[Proof of Theorem~\ref{thm:main}]
Fix one of the three families and work first over \(K=\Q(r)\).
Section~\ref{sec:translation} proves that the matrix \eqref{eq:familymatrix}
has rank at most four, including the specialization \(c_2=0\).
The appropriate nonzero value in \eqref{eq:witnesses} proves that the
four-row product matrix is invertible over \(K\). The four zero
polynomials in the certificate table are exactly
\eqref{eq:fourchecks}. Lemma~\ref{lem:finite} therefore gives
\eqref{eq:main} for every \(n\geq7\) as an identity in \(\Q(r)\).
Both sides are in \(\Z[r]\), so their difference is the zero polynomial
in \(\Z[r]\). This proves the denominator-free statement at every
parameter value. Division by \(D_k(r)\) where nonzero gives
\eqref{eq:normalized} and the proposed Somos--8 coefficients.
\end{proof}

\subsection{How the finite algebra is verified}

The primary verifier uses only the Python standard library. Integer
polynomials are stored as coefficient tuples in ascending degree. Addition
and multiplication are elementary operations, and every polynomial division
is checked for an exactly zero remainder.

Determinants are evaluated by fraction-free Bareiss elimination with row
pivoting. If \(p\) is the current pivot and \(p_{\rm old}\) the previous
pivot, an elimination update is
\begin{equation}
 a'_{ij}=\frac{p\,a_{ij}-a_{ik}a_{kj}}{p_{\rm old}}.
 \label{eq:bareiss}
\end{equation}
The previous pivot is initialized to one. Divisions in this algorithm are
exact; the implementation raises an exception if either an integer
remainder or a polynomial remainder occurs. A row exchange is accounted
for in the determinant sign. An identically zero pivot column returns a
zero determinant rather than dividing by zero.

The optional independent check constructs the moments afresh in SymPy's
\(\Z[r]\) domain and evaluates all the Hankel determinants using its
\texttt{DomainMatrix.det} implementation. In the execution accompanying
this report, all 21 moments and all 12 determinants, including the empty
one, agree for each family using SymPy 1.14.0.

A further exact audit evaluates the denominator-cleared recurrence for
all three families, each integer \(r=-5,-4,\ldots,7\), and each
\(n=7,8,\ldots,28\). All
\begin{equation}
 3\cdot13\cdot22=858
\end{equation}
checks pass, including roots of the denominators. These are useful tests
against implementation errors, but the all-index proof is the rank argument
plus the polynomial certificates, not the number 858.

\section{Singular parameters, zeros, and lower-order relations}
\label{sec:singular}

\subsection{Bilinear identity versus forward division}

Even when \(D_k(r)\ne0\), a forward update
\begin{equation}
 h_n=\frac{\sum_{j=1}^4P_{k,j}(r)h_{n-j}h_{n-8+j}}
              {D_k(r)h_{n-8}}
 \label{eq:forward}
\end{equation}
is only legal when \(h_{n-8}\ne0\). The polynomial statement
\eqref{eq:main} needs neither this assumption nor nonvanishing of other
Hankel determinants. The moment recursion and determinant definition always
select a sequence, even when a quotient recurrence does not uniquely
continue it.

For example, at \(r=0\), the families 11 and 12 have the same generating
function. Their Hankel determinants begin
\begin{equation}
 1,1,0,-1,-1,-1,0,1,1,1,0,-1,-1,-1,\ldots.
\end{equation}
The zero terms are already enough to show why unrestricted forward division
is not the right theorem. In this same specialization, family 11 yields
coefficients \((0,0,0,1)\), whereas family 12 yields
\((-5/3,0,5/3,1)\). Both identities hold: uniqueness of the coefficient
vector over \(\Q(r)\) does not imply uniqueness after a special rank drop.
We do not infer a global periodicity theorem here merely from the displayed
initial segment.

\subsection{Two explicit Somos--6 consequences of Conjecture 12}

At \(r=1\), the denominator \(D_{12}\) vanishes but the numerator vector is
\begin{equation}
 (P_{12,1},P_{12,2},P_{12,3},P_{12,4})=(-5,5,0,25).
\end{equation}
At \(r=3\), it is
\begin{equation}
 (-35,-105,980,735).
\end{equation}
After moving the outermost surviving product to the left and replacing
\(n-1\) by \(m\), Theorem~\ref{thm:main} gives the following.

\begin{corollary}[Order reduction at the exceptional parameters]
For the family \(k=12\), the following identities hold for every \(m\geq6\):
\begin{align}
 r=1:\qquad
 h_mh_{m-6}&=h_{m-1}h_{m-5}+5h_{m-3}^2,\label{eq:r1}\\
 r=3:\qquad
 h_mh_{m-6}&=-3h_{m-1}h_{m-5}
              +28h_{m-2}h_{m-4}+21h_{m-3}^2.
 \label{eq:r3}
\end{align}
\end{corollary}

These are genuine denominator-free consequences, not limits of an
ill-defined quotient. Their initial values provide useful direct examples:
\begin{center}
\begin{tabular}{c rrrrrrr}
\toprule
 &\(h_0\)&\(h_1\)&\(h_2\)&\(h_3\)&\(h_4\)&\(h_5\)&\(h_6\)\\
\midrule
\(r=1\)&1&1&$-2$&$-5$&$-13$&7&132\\
\(r=3\)&1&1&$-6$&$-49$&$-205$&$-1$&84864\\
\bottomrule
\end{tabular}
\end{center}
For instance, the first step of \eqref{eq:r3} reads
\(84864=3+34440+50421\).

The denominator roots of the other two families are also transparent from
\eqref{eq:D11} and \eqref{eq:D13}. We assert the specialized polynomial
identities there, but do not claim to classify every further reduction or
every vanishing numerator configuration.

\subsection{Two nonsingular integer examples}

For family 12 at \(r=2\), the recurrence becomes
\begin{equation}
 h_nh_{n-8}=5h_{n-1}h_{n-7}-18h_{n-2}h_{n-6}
             +77h_{n-3}h_{n-5}-13h_{n-4}^2.
 \label{eq:r2rec}
\end{equation}
The sequence, from \(h_0\), starts
\begin{equation}
 1,1,-4,-13,-43,87,2512,22041,179945,-6856935,
 -162649748,\ldots.
\end{equation}
At \(r=4\), the parameters are
\begin{equation}
 (\alpha,\beta,\gamma,\delta)
 =\left(-\frac{101}{3},-\frac{484}{3},4299,\frac{23359}{3}\right).
\end{equation}
This also proves the specific Somos--8 assertion in Barry's Example 14:
its denominator uses
\(x(1-3x)/(1-4x)=x+x^2/(1-4x)\), so its generating function is precisely
our family 12 at \(r=4\). The relevant moments and Hankel values can be
regenerated without using that assertion; files for \(r=0,1,2,3,4\) are
included in the package.

\section{A positive coefficient formula and a tree model}
\label{sec:combinatorics}

This section concerns the moments \(a_n^{(12)}\), not the determinant
sequence. In the second family, define
\begin{equation}
 L(x)=x+\frac{x^2}{1-rx}.
 \label{eq:L}
\end{equation}
Equation \eqref{eq:c12cf} is equivalent to
\begin{equation}
 g(x)=1+L(x)g(x)+x^3g(x)^2.
 \label{eq:tree}
\end{equation}
Thus for a nonnegative integer \(r\), \(a_n\) counts rooted ordered trees
with the following decorations. A leaf has size zero. A unary vertex is
of one of two types: a short unary vertex of size one, or a long unary
vertex of size two carrying an additional possibly empty word on an
\(r\)-letter alphabet, each letter contributing one to the size. A binary
vertex has two ordered children and size three. The size of a decorated
tree is the sum of its vertex and word sizes.

There are finitely many such trees of each size because every non-leaf
vertex has positive size. The root decomposition is exactly
\eqref{eq:tree}; hence this interpretation follows from formal uniqueness.
In particular \(a_n(r)\in\mathbb N[r]\), even though the compact recurrence
\eqref{eq:momentrec} contains subtraction.

For nonnegative integers \(j,\ell\), put
\begin{equation}
 W(j,\ell)=
 \begin{cases}
 1,&j=\ell=0,\\
 0,&j=0,\ \ell>0,\\
 \displaystyle\binom{j+\ell-1}{\ell},&j\geq1.
 \end{cases}
 \label{eq:W}
\end{equation}

\begin{proposition}[Finite positive sum]
\label{prop:coeff}
For every \(n\geq0\),
\begin{equation}
 a_n^{(12)}(r)=
 \sum_{\substack{k,q,j,\ell\geq0\\j\leq q\\3k+q+j+\ell=n}}
 C_k\binom{2k+q}{q}\binom qj W(j,\ell)r^\ell,
 \qquad C_k=\frac1{k+1}\binom{2k}{k}.
 \label{eq:positiveformula}
\end{equation}
The formula is a polynomial identity, so it also holds for negative or
complex \(r\), although the counting interpretation then no longer applies.
\end{proposition}
\begin{proof}
Solving \eqref{eq:tree} in Catalan form gives
\begin{equation}
 g(x)=\sum_{k\geq0}C_kx^{3k}(1-L(x))^{-2k-1}.
\end{equation}
Expand
\begin{equation}
 (1-L)^{-2k-1}=\sum_{q\geq0}\binom{2k+q}{q}L^q,
 \qquad
 L^q=\sum_{j=0}^q\binom qj x^{q+j}(1-rx)^{-j}.
\end{equation}
The coefficient of \(x^\ell\) in the last factor is
\(W(j,\ell)r^\ell\). Collecting all contributions to \(x^n\) gives
\eqref{eq:positiveformula}.

Alternatively, take a full binary tree with \(k\) binary vertices, in
\(C_k\) ways. It has \(2k+1\) vertices at which chains of unary vertices
can be inserted, including above its root. Distributing \(q\) unary
vertices among these positions gives \(\binom{2k+q}{q}\) choices.
Choose \(j\) of the \(q\) unary vertices to be long. Distribute a total of
\(\ell\) additional letters among their words in \(W(j,\ell)\) ways and
color the letters in \(r^\ell\) ways. The total size is
\(3k+q+j+\ell\), exactly as in the sum.
\end{proof}

The generating function also gives a manifestly nonnegative recursion:
\begin{equation}
 a_n=a_{n-1}+\sum_{j=0}^{n-2}r^ja_{n-2-j}
                    +\sum_{i=0}^{n-3}a_ia_{n-3-i},\qquad n\geq1,
 \label{eq:positiverec}
\end{equation}
with negative-index terms and empty sums zero. This provides a useful
independent way to check the compact moment recursion.

\section{Dominant-singularity asymptotics for the second family}
\label{sec:asymptotic}

Fix a real \(r\geq0\) throughout this section. Set
\begin{equation}
 \Delta(x)=(1-L(x))^2-4x^3,
 \label{eq:Delta}
\end{equation}
so that
\begin{equation}
 g(x)=\frac{1-L(x)-\sqrt{\Delta(x)}}{2x^3}.
 \label{eq:asymroot}
\end{equation}
We give the full inverse-power expansion associated with the dominant
singularity. This is not an asymptotic theorem for the Hankel determinants,
and it is not a claim to include exponentially smaller contributions from
every other algebraic singularity.

\begin{lemma}[Unique dominant singularity]
\label{lem:rho}
There is a unique positive number \(\rho\) satisfying
\begin{equation}
 \rho+\frac{\rho^2}{1-r\rho}+2\rho^{3/2}=1.
 \label{eq:rhoequality}
\end{equation}
It lies in \((0,1/r)\) when \(r>0\), and in \((0,1)\) when \(r=0\).
The function \(g\) has no singularity in \(|x|<\rho\), and \(\rho\) is
its only singularity on \(|x|=\rho\). Moreover,
\begin{equation}
 \Delta'(\rho)
 =-4\rho^{3/2}L'(\rho)-12\rho^2<0.
 \label{eq:Deltaprime}
\end{equation}
\end{lemma}
\begin{proof}
The left side of \eqref{eq:rhoequality} is strictly increasing from zero
on the stated positive interval and exceeds one before the upper endpoint.
This proves existence and uniqueness. The equality gives
\(1-L(\rho)=2\rho^{3/2}\), hence \(\Delta(\rho)=0\) and
\eqref{eq:Deltaprime}.

For \(|z|<\rho\), the nonnegative coefficients of \(L\) imply
\(L(|z|)+2|z|^{3/2}<1\). A zero of \(\Delta(z)\) would give a complex
number \(w\) with \(w^2=z^3\) and \(1=L(z)\pm2w\), contradicting
\[
 1\leq |L(z)|+2|z|^{3/2}\leq L(|z|)+2|z|^{3/2}<1.
\]
Thus the square-root branch extends analytically inside that disk, and its
apparent singularity at zero is removable.

If \(|z|=\rho\) and \(\Delta(z)=0\), equality must hold throughout the
same triangle inequalities. In particular \(|L(z)|=L(\rho)\). The
coefficients of both \(x\) and \(x^2\) in \(L\) are positive; equality
forces these two summands to have the same argument. Writing
\(z=\rho e^{i\theta}\), this requires
\(e^{i\theta}=e^{2i\theta}\), so \(z=\rho\).
The rational coefficients have no pole on the disk because \(r\rho<1\).
Finally, the simple zero \eqref{eq:Deltaprime} gives a nonzero square-root
term, so \(\rho\) really is a singularity.
\end{proof}

\begin{theorem}[Leading term and first correction]
\label{thm:asymptotic}
Define
\begin{equation}
 \lambda=\frac{\sqrt{-\rho\Delta'(\rho)}}{2\rho^3},\qquad
 C=\frac{\lambda}{2\sqrt\pi},\qquad
 c_1=-\frac{33}{8}+\frac{3\rho\Delta''(\rho)}{8\Delta'(\rho)}.
 \label{eq:asymconstants}
\end{equation}
Then, as \(n\to\infty\),
\begin{equation}
 a_n^{(12)}(r)=C\rho^{-n}n^{-3/2}
             \left(1+\frac{c_1}{n}+O(n^{-2})\right).
 \label{eq:asymptotic}
\end{equation}
The estimate is for each fixed \(r\geq0\); no uniformity in an unbounded
range of \(r\) is asserted.
\end{theorem}
\begin{proof}
Set \(t=1-x/\rho\). The nonanalytic part of \eqref{eq:asymroot} is
\begin{equation}
 -\lambda t^{1/2}R(t),\qquad
 R(t)=(1-t)^{-3}
 \sqrt{\frac{\Delta(\rho(1-t))}{-\rho\Delta'(\rho)t}}.
 \label{eq:R}
\end{equation}
The quotient under the square root is analytic at zero with value one.
Therefore \(R\) is analytic there, \(R(0)=1\), and
\begin{equation}
 R(t)=1+R_1t+O(t^2),\qquad
 R_1=3-\frac{\rho\Delta''(\rho)}{4\Delta'(\rho)}.
 \label{eq:R1}
\end{equation}
The algebraic function has a suitable indented complex neighborhood of its
unique dominant singularity, by Lemma~\ref{lem:rho}; the other zeros and
poles are separated from it. Singularity transfer \cite{FlajoletOdlyzko}
applies to these half-integer terms.

For nonintegral \(a\), the binomial coefficient identity is
\begin{equation}
 [x^n](1-x/\rho)^a
 =\rho^{-n}\frac{\Gamma(n-a)}{\Gamma(-a)\Gamma(n+1)}.
 \label{eq:gammaidentity}
\end{equation}
For \(a=1/2\), its gamma ratio is
\(n^{-3/2}(1+3/(8n)+O(n^{-2}))\), and
\(\Gamma(-1/2)=-2\sqrt\pi\). For \(a=3/2\), the relative contribution
to the first correction is \(-3R_1/(2n)\). Hence the correction is
\(3/8-3R_1/2\), which simplifies to \eqref{eq:asymconstants}.
The analytic part does not contribute to these algebraic-order dominant
terms. This proves \eqref{eq:asymptotic}.
\end{proof}

\subsection{An explicit all-orders prescription}

Write \(R(t)=\sum_{j\geq0}R_jt^j\) in \eqref{eq:R}. For every fixed
integer \(J\geq0\), singularity transfer gives
\begin{align}
 a_n={}&-\lambda\rho^{-n}
 \sum_{j=0}^{J}R_j
 \frac{\Gamma(n-j-\tfrac12)}{\Gamma(-j-\tfrac12)\Gamma(n+1)}
 \notag\\
 &\hspace{2em}{}+O\!\left(\rho^{-n}n^{-J-5/2}\right).
 \label{eq:allorders}
\end{align}
Thus the coefficients of the complete dominant inverse-power expansion are
obtained by differentiating the explicit analytic function \(R\) and
expanding elementary gamma ratios. This prescription avoids an auxiliary
unknown recurrence for the asymptotic coefficients. It also makes the
error order at every finite truncation explicit.

At \(r=2\), the computations supplied with this report give
\begin{align}
 \rho&\approx0.32542163481800348746,\notag\\
 \rho^{-1}&\approx3.07293644001039978597,\notag\\
 C&\approx5.05502055022293177233,\notag\\
 c_1&\approx-4.11576700299955384839.
 \label{eq:numericalconstants}
\end{align}
The table compares exact coefficients with the first two approximations:
\begin{center}
\begin{tabular}{r rr}
\toprule
\(n\)&\(a_n/(C\rho^{-n}n^{-3/2})\)
 &\(a_n/(C\rho^{-n}n^{-3/2}(1+c_1/n))\)\\
\midrule
25 &0.8421704141&1.0081414221\\
50 &0.9207369642&1.0033260927\\
100&0.9597064447&1.0009012063\\
200&0.9796476789&1.0002312733\\
400&0.9897684471&1.0000584662\\
\bottomrule
\end{tabular}
\end{center}
The numerical script uses 80-digit decimal arithmetic for the constants
and integer arithmetic for the coefficients. These values illustrate the
proved asymptotics; they are not used in the Somos proof.

\section{Reproducibility and audit trail}
\label{sec:reproduce}

The archive is organized so that the exact proof checks can be run without
installing a computer algebra system. From its root directory, execute
\begin{lstlisting}
python3 code/verify_certificates.py --audit
python3 code/examples.py
\end{lstlisting}
The first command recomputes all certificate polynomials and the rank
witnesses, then performs the 858 integer audits. The second regenerates
sequence files and the asymptotic comparison data. An optional independent
symbolic check is
\begin{lstlisting}
python3 code/cross_check_sympy.py
\end{lstlisting}
which requires SymPy. The standard-library verifier requires Python 3.10
or later and makes no network calls.

The main certificate file records integer coefficients in ascending powers
of \(r\). For example, \([1,-2,3]\) denotes \(1-2r+3r^2\), and the zero
polynomial is \([0]\). Its moment array is indexed from 0 through 20; its
Hankel array is indexed from \(-1\) through 10. This convention is explicitly
stored and used by both verifiers, avoiding silent off-by-one shifts.

The delivered files have the following roles:
\begin{center}
\small
\begin{tabular}{p{0.43\textwidth}p{0.49\textwidth}}
\toprule
File or directory&Purpose\\
\midrule
\texttt{article.tex}, \texttt{article.pdf}&The mathematical report and its editable source.\\
\texttt{code/verify\_certificates.py}&Exact integer-polynomial determinant proof checks.\\
\texttt{code/cross\_check\_sympy.py}&Independent optional symbolic audit.\\
\texttt{code/examples.py}&Exact moment/Hankel data and decimal asymptotic comparisons.\\
\texttt{data/certificates.json}&All finite polynomial certificates and rank witnesses.\\
\texttt{data/numeric\_audit.json}&The scope and result of the 858 auxiliary checks.\\
\texttt{data/C12\_r*\_*.txt}&Indexed moment and Hankel sequences at five integer parameters.\\
\texttt{data/asymptotics\_C12\_r2.json}&High-precision constants and comparison ratios.\\
\texttt{README.md}, \texttt{build.sh}&Usage, proof dependencies, and compilation instructions.\\
\bottomrule
\end{tabular}
\end{center}

The TeX source compiles with pdfLaTeX. The build script runs it twice to
resolve references and the table of contents. No font files or copies of
third-party papers are redistributed. The bibliography links to the
original sources.

\section{Scope, significance, and remaining questions}
\label{sec:conclusion}

The main outcome is a proof of three precise published Somos--8 assertions,
not merely the discovery of matching initial terms. The decisive observation
is that the universal genus-two theorem should be used as a \emph{polynomial
rank bound}, rather than by substituting into its potentially singular
rational coordinate formulas. This both validates the degenerate
\(c_2=0\) specialization and makes a small family-specific certificate
sufficient.

The work also separates three issues that can otherwise be conflated.
Integrality follows from the moment construction. A bilinear identity follows
from the rank theorem and exact coefficient identification. A division-based
recurrence additionally requires nonzero divisors. Our theorem settles the
first two without silently assuming the third.

The exceptional values \(r=1,3\) in the second family are not discarded;
they yield explicit order-six relations. The positive tree formula and the
asymptotic expansion add information about the underlying moment sequence.
They do not provide Hankel determinant asymptotics, which would require
further analysis of the relevant genus-two dynamics and its degenerations.

Several natural directions remain: classifying all rank drops and all zero
patterns at specialized parameters; obtaining similarly compact certificates
for other low-degree algebraic moment families; developing a direct
family-specific determinant or continued-fraction proof that does not import
Hone's general theorem; and deriving asymptotics for the determinant sequences
themselves. These directions are not asserted to be solved here.

Finally, the novelty qualification remains important. The theorem and its
certificates are fully specified and independently reproducible, but an
unlocated prior application of Hone's theorem would change the historical
claim, not the validity of these identities. External mathematical review
would be appropriate before treating the report as a publication-ready
priority claim.

\appendix
\section{The other two rank witnesses}
\label{app:witness}

For reference, the matrices formed from rows \(7,8,9,10\) at \(r=2\)
for the other two families are
\begin{equation}
 M_{11}(2)=
 \begin{pmatrix}
 -135120&-34593&-1430&1681\\
 14703657&-270240&472771&511225\\
 1376443582&-200949979&96610800&132963961\\
 -2390410472409&-492078580565&-56515956289&18257414400
 \end{pmatrix},
\end{equation}
and
\begin{equation}
 M_{13}(2)=
 \begin{pmatrix}
 -1733&-1186&-78&169\\
 57264&-1733&7709&6084\\
 1058325&-372216&135174&351649\\
 -259181611&-82549350&-16978776&3003289
 \end{pmatrix}.
\end{equation}
Their exact determinants are listed in \eqref{eq:witnesses}. The witness
calculation does not claim that \(M_k(r)\) is invertible for every \(r\).
It proves invertibility over the rational function field, which is all the
finite-certification lemma needs before polynomial specialization.

\section{A compact logical checklist}
\label{app:checklist}

The proof can be audited in the following order.
\begin{enumerate}[leftmargin=2em,itemsep=3pt]
\item Confirm the three continued fractions and recurrence parameters in
Barry's Conjectures 11--13, with the stated correction to the separate
Catalan display for Conjecture 13.
\item Verify the compact moment recursion \eqref{eq:momentrec} by multiplying
out \eqref{eq:uniform}.
\item Verify the universal moment form \eqref{eq:universal} and its match to
Hone's Theorem 5.4 via \eqref{eq:PGST}.
\item Pass to polynomial minors before setting \(c_2=0\), using
Lemma~\ref{lem:specialization}.
\item Check the binomial translation and the seven substituted parameters
in \eqref{eq:sevenparameters}.
\item Recompute the twelve zero residual polynomials and the three nonzero
rank witnesses.
\item Apply Lemma~\ref{lem:finite}, then specialize the resulting
\(\Z[r]\)-identities, without dividing by a Hankel determinant.
\end{enumerate}
Every family-specific claim needed for the all-index result is covered by
these steps. The auxiliary numerical tests and the asymptotic section are
not premises of the main theorem.

\begin{thebibliography}{9}
\addcontentsline{toc}{section}{References}

\bibitem{Barry}
P. Barry,
\emph{Conjectures on Somos 4, 6 and 8 sequences using Riordan arrays and the
Catalan numbers}, arXiv:2211.12637v1, 22 November 2022.
Conjectures 11--13 are in Section 6, pp.\ 8--10; Example 14 is on p.\ 10.
\url{https://arxiv.org/abs/2211.12637}.

\bibitem{Hone}
A. N. W. Hone,
\emph{Continued fractions and Hankel determinants from hyperelliptic curves},
Communications on Pure and Applied Mathematics \textbf{74} (2021),
2310--2347, doi:10.1002/cpa.21923.
The version used for equation and theorem numbers is arXiv:1907.05204v3,
30 December 2019, especially equation (4.22) and Theorem 5.4.
\url{https://arxiv.org/abs/1907.05204}.

\bibitem{WangZhang}
Y. Wang and Z. Zhang,
\emph{A sufficient condition for $(\alpha,\beta)$ Somos 4 Hankel determinants},
arXiv:2305.05995v2, 15 June 2023; published in Discrete Mathematics (2024).
The four corollaries address Barry's Somos--4 conjectures.
\url{https://arxiv.org/abs/2305.05995}.

\bibitem{LiuWangZhang}
F. Liu, Y. Wang, and Z. Zhang,
\emph{Proof of Barry's Four Hankel Determinant Conjectures},
arXiv:2607.18644, 2026. This paper concerns the central transform of four
rational families, a distinct set of conjectures.
\url{https://arxiv.org/abs/2607.18644}.

\bibitem{Scheuerle}
T. Scheuerle,
\emph{Direct Generation of a Somos--4 Sequence from an Algebraic Generating
Function}, arXiv:2609.15754v1, 14 September 2026.
\url{https://arxiv.org/abs/2609.15754}.

\bibitem{FlajoletOdlyzko}
P. Flajolet and A. Odlyzko,
\emph{Singularity analysis of generating functions},
SIAM Journal on Discrete Mathematics \textbf{3} (1990), 216--240,
doi:10.1137/0403019.
\url{https://algo.inria.fr/flajolet/Publications/FlOd90b.pdf}.

\end{thebibliography}
\end{document}
